(1) 占据数 \(f(\epsilon)=1/(z^{-1}e^{\beta\epsilon}-\eta)=\sum_{\ell\ge1}\eta^{\ell-1}z^\ell e^{-\ell\beta\epsilon}\)。对三维动量高斯积分，得到
\[
n\lambda_T^3=z+\eta z^2/2^{3/2}+O(z^3),\quad
P\lambda_T^3/(k_BT)=z+\eta z^2/2^{5/2}+O(z^3).
\]
压力公式也可由 \(\ln\Xi=(V/\lambda_T^3)\sum_{\ell\ge1}\eta^{\ell-1}z^\ell/\ell^{5/2}\) 求得。
(2) 设 \(x=n\lambda_T^3\)，反演 \(z=x-\eta x^2/2^{3/2}+O(x^3)\)。代回压力得
\[
P=nk_BT[1-\eta x/2^{5/2}+O(x^2)],\quad B_2=-\eta\lambda_T^3/2^{5/2}.
\]
玻色交换降低压强，费米排斥提高压强；这里无真实粒子间相互作用。要求 \(x\ll1\)，不仅仅是宏观密度小。
(3) \(\mu=k_BT\ln z=k_BT[\ln x-\eta x/2^{3/2}+O(x^2)]\)。用 \(F=\mu N-PV\) 得
\[
\frac F{Nk_BT}=\ln x-1-\eta\frac{x}{2^{5/2}}+O(x^2).
\]
在固定 \(T\) 下 \(x\propto n\)，所以 \(n^2\partial(F/N)/\partial n=nk_BT[1-\eta x/2^{5/2}]\)，与第二问独立吻合。
(4) 自由粒子动能满足 \(U=3PV/2\)，因此 \(U/N=\frac32k_BT[1-\eta x/2^{5/2}]+O(x^2)\)。熵由 \(S=(U+PV-\mu N)/T\) 得
\[
\frac S{Nk_B}=\frac52-\ln x-\eta\frac{x}{2^{7/2}}+O(x^2).
\]
固定 \(n\) 时 \(x\propto T^{-3/2}\)，故高温量子修正消失，恢复Sackur--Tetrode形式（无内部简并）。低温或高密使 \(x\gtrsim1\)，玻色凝聚或费米简并须用完整分布，截断维里级数无效。
